% =====================================================================
%  ACRÓNIMOS
%  Define aquí cada sigla UNA vez. En el texto usa \ac{API}: la primera vez
%  sale «Application Programming Interface (API)» y después solo «API».
%  Gracias a [printonlyused], la lista solo muestra las siglas que usas.
%  Otras formas: \acs{API} (solo sigla), \acl{API} (solo forma larga),
%  \acp{API} (plural).
% =====================================================================
\cleardoublepage
\chapter*{\acronimosname}
\pdfbookmark[0]{\acronimosname}{acronimos}

\begin{acronym}[RGPD]  % entre corchetes, la sigla más larga (para alinear)
  \acro{API}{\emph{Application Programming Interface}}
  \acro{EPS}{Escola Politècnica Superior}
  \acro{RF}{Requisito Funcional}
  \acro{RGPD}{Reglamento General de Protección de Datos}
  \acro{RNF}{Requisito No Funcional}
  \acro{REST}{\emph{Representational State Transfer}}
  \acro{SUS}{\emph{System Usability Scale}}
  \acro{TFG}{Trabajo de Fin de Grado}
  \acro{TFM}{Trabajo de Fin de Máster}
  \acro{UdL}{Universitat de Lleida}
\end{acronym}
